\begin{corollary}[Positive-rung well-definedness]
Let \(G=(V,E)\) be a finite loopless multigraph with individual edge
identities, let \(e_0\in E\), and let the palette have \(k\) colors.  Let
\(\mathcal E\) be a positive-rung extended Tashkinov tree with rung number
\(n\), input coloring \(\varphi\), certified series \(\mathcal T\), and
displayed tree exactly the final tree \(T_{n+1}\) of \(\mathcal T\).  Suppose
that \(\mathcal E\) has \noderef{MaximumProperty},
\(\noderef{MaximumDegree}(G)+1\leq k\), and
\(11\leq\noderef{TashkinovOrder}(G)\).

Then one of the following two alternatives holds.
\begin{enumerate}
\item There are an input coloring \(\varphi^*\), a certified series
      \(\mathcal T^*\) with exactly \(n\) noninitial iterations, and an ETT
      \(\mathcal E^*\) displayed on exactly the final tree of
      \(\mathcal T^*\), with generating coloring exactly its terminal
      coloring, such that the rung of \(\mathcal E^*\) is exactly \(n\),
      its dependent series is exactly \(\mathcal T^*\), its terminal
      coloring is proper and leaves exactly \(e_0\) uncolored, its final
      vertex set is both closed and strongly closed, and
      \(\mathcal E^*\) has both \noderef{MaximumProperty} and
      \noderef{TerminalClosureMaximum}.
\item There are an input coloring \(\varphi^*\), a certified series
      \(\mathcal T^+\) with exactly \(n+1\) noninitial iterations, and an
      ETT \(\mathcal E^+\) displayed on exactly the final tree of
      \(\mathcal T^+\), with generating coloring exactly its terminal
      coloring, such that the rung of \(\mathcal E^+\) is exactly \(n+1\),
      its dependent series is exactly \(\mathcal T^+\), its terminal
      coloring is proper and leaves exactly \(e_0\) uncolored, its final
      vertex set is closed, and \(\mathcal E^+\) has
      \noderef{MaximumProperty}.
\end{enumerate}
No correspondence assertion between the original ETT and the ETT in the
second alternative is made.
\end{corollary}
\begin{proof}
Apply \noderef{NormalizePositiveTerminalMaximum} to \(\mathcal E\).  Its
hypotheses are exactly the assumed final-tree equality, maximum property,
degree bound, and large-order bound.  It supplies a proper coloring
\(\sigma_n\) leaving exactly \(e_0\) uncolored, a seed \(S\), a maximum
closure \(T^*\), and a
\noderef{CorrespondingExtendedTashkinovTree} certificate \(C\).  It also
supplies the following statement-level conclusions: \(T^*\) is the exact
final tree of \(C\)'s replacement series; the replacement ETT has
\noderef{MaximumProperty} and \noderef{TerminalClosureMaximum}; and
\(T^*\) is a \noderef{TashkinovClosure} under \(\sigma_n\), hence its
vertex set is closed.

Unpack the data in \(C\).  By the definition
\noderef{CorrespondingExtendedTashkinovTree}, there are an actual
replacement input coloring \(\varphi^*\), an actual certified replacement
series \(\mathcal T^*\) of the same rung \(n\), and an actual replacement
ETT whose dependent series is \(\mathcal T^*\) and whose final generating
coloring is \(\sigma_n\).  Eliminate the rung equality first and then the
dependent equality of certified series.  Next eliminate the equality of
the generating coloring with \(\sigma_n\), and finally the equality of
the replacement series' exact final tree with \(T^*\).  After these four
dependent transports, the replacement ETT may be presented as an ETT
\(\mathcal E^*\) whose displayed tree is literally
\(\mathcal T^*.\operatorname{tree}(\operatorname{last}(n+1))\), whose
generating coloring is literally
\(\mathcal T^*.\operatorname{coloring}(\operatorname{last}(n))\), whose
rung is \(n\), and whose dependent series is \(\mathcal T^*\).  The maximum
property and terminal maximum property returned by
\noderef{NormalizePositiveTerminalMaximum} transport along these same
equalities to \(\mathcal E^*\).  The properness and exactly-one-uncolored
conclusions for \(\sigma_n\) become those properties for the exact terminal
coloring of \(\mathcal T^*\), while the closure conclusion for \(T^*\)
becomes closedness of the exact final vertex set.

Test whether that exact final vertex set is strongly closed under the exact
terminal coloring.  If it is, retain this fact together with the rung and
dependent-series equalities, properness, exactly-one-uncoloredness,
closedness, maximum property, and terminal maximum property just obtained.
These witnesses give the first alternative.

Suppose instead that the exact final vertex set is not strongly closed.
Apply \noderef{AlgorithmThreeOneStepExists} to the certified replacement
series \(\mathcal T^*\).  The required degree bound is assumed; the required
terminal properness, exactly-one-uncoloredness, and closedness were obtained
above; and the last hypothesis is the present failure of strong closedness.
The theorem returns a certified series \(\mathcal T^+\) with \(n+1\)
noninitial iterations whose complete data are exactly the data of
\(\mathcal T^*\) followed by the supplied Algorithm~3.1 step.  In
particular, this is the exact append equation required by
\noderef{MaximumPropertyAppend}.  Applying that lemma to
\(\mathcal E^*\), its maximum property, its terminal maximum property, and
this exact append produces an exact-final ETT \(\mathcal E^+\).  The same
lemma states, rather than merely choosing definitionally convenient data,
that \(\mathcal E^+\) has rung \(n+1\), has dependent series exactly
\(\mathcal T^+\), and has \noderef{MaximumProperty}.

It remains only to verify the three terminal coloring facts demanded in the
second alternative.  Put \(r=n+1>0\) and \(q=r-1\).  Then the successor of
\(q\) is the final coloring index and its second successor is the final tree
index.  By \noderef{TashkinovSeries}, the exact last index of
\(\mathcal T^+\) has a certificate classified by
\noderef{TashkinovSeriesStep}.  The initial constructor is impossible,
because it certifies index zero whereas this final index has positive value
\(r\).  There are three remaining cases.

In the SE case, the constructor identifies the terminal coloring with the
selected maximum-defective coloring.  The fields of
\noderef{MaximumDefectiveChoice} state that this coloring is proper and
leaves exactly \(e_0\) uncolored, and the closed-extension field states that
the final tree is a closure under that coloring.  Thus its final vertex set
is closed.  In the PE case, the constructor directly records terminal
properness and exactly-one-uncoloredness, and its closed-extension field
again records closedness of the exact final tree.  In the RE case, the
constructor identifies the terminal coloring with the preceding proper
coloring that leaves exactly \(e_0\) uncolored, and its closed-extension
field supplies closedness.  Rewriting the successor indices \(q+1\) and
\(q+2\) as the exact final coloring and tree indices gives, in all three
cases, properness, exactly-one-uncoloredness, and closedness for the terminal
data of \(\mathcal T^+\).

Combine these three terminal facts with the exact rung equality, dependent
series equality, and maximum property supplied by
\noderef{MaximumPropertyAppend}.  This gives the second alternative.  The
only correspondence certificate used was the same-rung certificate needed
to normalize \(\mathcal T^*\); after the Algorithm~3.1 append we assert only
the exact new series and ETT data furnished by the append lemma, so no
same-rung correspondence is claimed across the append.
\end{proof}
